\scanpage{82}
% loc 1
de la première des relations
% loc 2–3
\[ \left\{\begin{aligned}\operatorname{Tr}1&=2\\ \det 1&=1\end{aligned}\right. \tag*{(5)} \]
% loc 4–6
\deleted{les quelles}\inserted{Dont la 2°} se déduit \deleted{\unsure{est}} de $(2')$ en y faisant $u=1$ (et $\det 1=1$ est donné comme connu).

% loc 7–8
En polarisant (2), on trouve pour $u,v\in A$ (tenant compte de 3)
% loc 9
\[ uv+vu-\operatorname{Tr}(\unsure{v|u})u-\operatorname{Tr}(u)\unsure{v|u}+(\operatorname{Tr}u\operatorname{Tr}v-\operatorname{Tr}(uv))1 \tag*{(6)} \]
% loc 10
soit
% loc 11
\[ uv+vu=\operatorname{Tr}v)u+(\operatorname{Tr}u)v+(\operatorname{Tr}(uv)-\operatorname{Tr}u\operatorname{Tr}v)1 \tag*{$(6')$} \]
% loc 11–20, oblique note from left edge
% slanted
\margin{\uline{NB} Cette relation $(6')$ (pour la seule donnée de la forme $\operatorname{Tr}$, sans dét ni \unsure{Ham. Cayley|Ham.-Cayley}.) équivaut \unsure{``au} fait que $u\mapsto\bar u'=\operatorname{Tr}(u).1-u$ est un ``antiautomorphisme'' (\unsure{en vertu de|en vertu} \deleted{de carré l'\unsure{id.} identité}). Ceci dit, pour les \unsure{deux} ``formes'' $\operatorname{Tr}$, \unsure{l'une|la} quadratique $u\mapsto\det u$, les deux relations de $(2)\ (3)\ (6)$ ($\unsure{\Longleftrightarrow|\Rightarrow}6''$) impliquent la troisième.}
% loc 12–13
\deleted{\uline{Corollaire} On en déduit aussitôt} Cela \deleted{signif} implique que mod
% loc 14
\[ k.1+ku+kv, \]
% loc 15
$uv$ et $vu$ sont \deleted{\unsure{opposés}} opposés
% loc 16
\[ vu\equiv uv\qquad\bmod\quad k1+ku+kv. \]
% loc 17–20, proposition
\uline{Proposition} Soit $A$ avec (1) (2) (3) (d'où $(6')$) Alors $\forall u,v\in A$, \deleted{$k+ku+kv+k.\unsure{v|uv}$} le sous $k$-module engendré par $1,u,v,uv$ est une sous-Algèbre.

% loc 21, bottom
Il suffit de voir que \deleted{le produit de}
% archive sheet number: 81
